\chapter{后缀数组}\label{ch:suffix-array}

\chaptermeta{suffix_array.hpp}{def.hpp}{$O(n\log n)$，LCP $O(n)$}

对任意可随机访问、元素可比较的序列构造后缀数组，给出字典序排名、每个起点的排名以及相邻后缀的最长公共前缀。

\section{接口}

\apitable{
    \api{SuffixArray(const R& s)}{构造；\code{R} 需支持 \code{size()} 与随机访问，元素之间可比较大小}
    \api{usize size()}{序列长度 $n$}
    \api{usize sa(usize i)}{字典序第 $i$ 小的后缀的起始位置（$0\le i<n$）}
    \api{usize rk(usize i)}{从位置 $i$ 开始的后缀的排名，是 \code{sa} 的逆映射}
    \api{usize height(usize i)}{第 $i$ 小与第 $i-1$ 小后缀的 LCP；\code{height(0) == 0}}
}

\section{约定}

元素只需支持 \code{<} 与 \code{!=}，不必是字符；整数序列、\code{std::string} 与 \code{std::vector<int>} 都可直接传入。空序列合法，但此时不要调用 \code{sa}、\code{rk}、\code{height}。\code{sa} 与 \code{rk} 互为逆映射，下标均从 0 开始；\code{height(0)} 没有「前一个后缀」，按定义为 0，求任意两后缀的 LCP 需要另做 RMQ。

\fullsource{suffix_array.hpp}
